\chapter[Introduction]{Introduction}
Wildlife monitoring is a key information source for ecologists and biologists, underpinning a broad array of research in wildlife ecology. It aids in revealing the complex dynamics within ecosystems, assessing the effects of human activities and climate change, and informing effective species conservation measures~\citep{ahumada2013monitoring, he2016visual, tuia2022perspectives, bruce2025large, lindenmayer2022we}. Data collection for these purposes can employ a range of sensors~\citep{tuia2022perspectives}, providing multiple perspectives on animal populations. Examples include remote sensing (satellite or drone images, GPS tags), fixed devices (camera traps and bioacoustic recorders), and citizen-science platforms where members of the public submit observations via systems such as iNaturalist~\citep{horn2018inat, shepley2021automated, herrera2025inaturalist} and eButterfly~\citep{prudic2017ebutterfly}.

\section{Context and motivation}

Camera traps are established, non-invasive, and cost-effective tools for recording wildlife in situ. The data they produce support ecological analyses of terrestrial vertebrate communities, including studies of species occupancy, diversity, abundance, density, community composition, and behavior~\citep{ahumada2013monitoring, beery2021iwildcam, tanwar2021camera, delisle2021next}. These devices are typically triggered by motion or programmed to capture images at scheduled intervals, producing short image or video sequences~\citep{he2016visual}. As a result, camera-trap deployments often generate datasets at a scale that is difficult to process manually. For example, the Snapshot Serengeti project collected 3.2 million images between June 2010 and May 2013 across six seasons~\citep{swanson2015snapshot}, and this collection has since grown to 7.1 million images over eleven seasons~\citep{ss_lila}. Similarly, the Caltech Camera Traps dataset contains 243,100 images from 140 camera locations in the Southwestern United States, annotated across 21 animal categories plus empty images, with approximately 66,000 bounding box annotations. Around 70\% of its images are labeled as empty~\citep{beery2018recognition}. The annotation burden grows with this scale. The original Snapshot Serengeti labeling effort required more than 17,000 volunteer hours~\citep{norouzzadeh2018automatically}, and mobilizing comparable effort is often impractical for smaller studies or for projects focused on less charismatic species.

Over the last decade, deep learning has become a central approach for image analysis and has been widely applied to automate camera-trap data processing. Representative applications include empty-image filtering~\citep{mononen2025active, yousif2019animal, yang2024systematic, cunha2021filtering}, species classification~\citep{chen2023applying, velez2023evaluation, beery2018recognition, norouzzadeh2018automatically}, individual counting~\citep{norouzzadeh2018automatically, norouzzadeh2021deep, arshad2020my, wang2022deep}, and behavior recognition~\citep{schindler2021identification, bain2021automated, binta2023animal, lu2023mammalclub}. In parallel, several software platforms have emerged to help researchers accelerate labeling and information extraction with deep-learning models~\citep{velez2023evaluation}. Some provide user-friendly interfaces that allow a wider community of camera-trap users to train custom models. For example, MegaDetector~\citep{beery2019efficient} is a general detector for animals, people, and vehicles that helps identify relevant frames. The MLWIC2 package provides pre-trained models and training tools for species classification and empty-frame filtering in R~\citep{tabak2020improving}. Other GUI-based platforms, such as Wildlife Insights, Conservation AI, IFLA, and Animal Scanner, have also been developed to simplify information extraction from camera-trap images~\citep{ahumada2020wildlife, chalmers2019conservation, xi2021image, yousif2019animal}.

These AI-driven tools have increased the speed and scalability of camera-trap data processing and annotation, enabling applications at larger scales~\citep{fennell2022use, clarfeld2023evaluating, bohner2023semi, lonsinger2024efficacy}. Recent studies further show the use of these tools in broader and more diverse ecological settings~\citep{celis2024versatile, chen2024minimizing}. Species labels produced by computer-vision models can approximate expert-derived ecological metrics, such as species richness, occupancy, and activity patterns~\citep{whytock2021robust, pantazis2024deep}. Nonetheless, important obstacles persist, including class imbalance and limited generalization across domains~\citep{beery2021scaling, yang2021systematic, zhu2022class, malik2021two}. These limitations are also reported in studies of cross-domain camera-trap recognition~\citep{cunha2023bag}. Moreover, a large portion of camera-trap data still relies on manual processing, which limits the scalability of biodiversity monitoring because of the human labor involved~\citep{beery2023wild}.

\section{Challenges in computer vision for camera-trap data analysis}

Extracting relevant information from camera-trap images involves several difficulties that can hinder accurate analysis, even for experienced specialists. Common problems reported in the literature include animals appearing far from the camera, natural camouflage, complex or unusual body postures, and inadequate lighting conditions~\citep{villa2017towards, norouzzadeh2018automatically, beery2018recognition, gomez2016animal}. Under such circumstances, capturing multiple images per trigger event can improve identification accuracy, as different poses may reveal distinctive morphological features needed for species recognition~\citep{beery2018recognition}.

Domain generalization is another key challenge. Because camera traps are stationary, models can inadvertently learn correlations with site-specific backgrounds present in the training data. This can reduce performance when models are deployed in new environments~\citep{beery2018recognition, schneider2020three, tabak2019machine, norman2023can}.

Beyond these aspects, camera-trap datasets often contain additional complexities, including annotation errors, temporal and spatial dependencies among observations, and limited dataset sizes~\citep{norouzzadeh2018automatically, schneider2020three, brookes2025panaf, bevan2025deep}. Collectively, these challenges motivate methodological approaches that are better aligned with the characteristics of wildlife monitoring data.

\section{Computer vision tasks in automated analysis of camera-trap data}

Before ecologists can incorporate camera-trap data into ecological models, raw images must be processed to extract the information required by the study. This information may include which species are present, how many individuals of each species appear in a capture event, the sex or age class of the animals, and their observed behaviors, such as feeding, resting, or fighting~\citep{norouzzadeh2018automatically}. The required outputs depend on the research question or project goals, and these processing steps can be formulated as computer-vision problems that motivate tailored methods.

A widely investigated task is species recognition, usually formulated as a visual classification problem in which each image is treated as an independent example. Although classification is a classical problem in computer vision, camera-trap datasets introduce several difficulties, including class imbalance~\citep{schneider2020three, johnson2019survey, li2024deep}, few-shot learning scenarios~\citep{alencar2024zero}, and fine-grained distinctions between visually similar species~\citep{van2017devil}. When species-level identification is not possible, labels at higher taxonomic ranks can still have ecological value~\citep{schneider2024recognition}, which introduces an additional hierarchical dimension to the classification problem. Related objectives, such as estimating age, determining sex, and recognizing activities, can also be formulated as classification tasks~\citep{norouzzadeh2018automatically}.

Because many deployments capture multiple images per trigger event to increase the chance of obtaining an informative view, another important problem is species identification using all images from a single capture event~\citep{beery2018recognition, sundaresan2025adapting}. Standard workflows often treat each photo independently during both training and inference. However, exploiting the full image sequence can improve classification when discriminative markings or poses appear only in some frames. In this setting, the goal is to infer species identity from the complete sequence associated with a capture event rather than from a single image.

Given that many camera-trap frames are empty because of incidental sensor activations, automating the removal of empty images is also a practical task. Tools such as MegaDetector have been widely used for this purpose~\citep{beery2019efficient, gadot2024crop, velez2023evaluation, leorna2022human}. Empty-image filtering reduces the amount of data that must be reviewed manually and can serve as an initial step before downstream tasks such as species classification and behavior recognition.

\section{Research Hypotheses}

This thesis is guided by three hypotheses concerning the use of multimodal pre-trained models in camera-trap workflows under limited supervision.

\begin{description}
  \item[H1] Pre-trained multimodal models can support core camera-trap tasks, including empty-image filtering, species classification, and behavior identification, under zero-shot and few-shot settings without requiring task-specific training for each task.

  \item[H2] In zero-shot species classification, combining visual similarity with textual similarity can improve top-1 accuracy relative to a visual-only baseline, although the magnitude of the improvement may depend on the dataset and image setting.

  \item[H3] A low-supervision multimodal approach can provide empirical evidence for reducing dependence on annotated camera-trap data while identifying the conditions under which multimodal models and visual-textual fusion are most useful for biodiversity monitoring workflows.
\end{description}

\section{Objectives}

\subsection{General Objective}

To investigate the conditions under which multimodal machine learning models can support and improve the automation of camera-trap workflows for biodiversity monitoring, with emphasis on ecological vision tasks under low-supervision and zero-shot settings.

\subsection{Specific Objectives}

\begin{description}

\item[1.] To evaluate the potential of multimodal foundation models for core camera-trap tasks, including empty-image filtering, species classification, and behavior identification.

\item[2.] To assess how the performance of multimodal models varies across task types, supervision regimes, and input contexts in camera-trap data.

\item[3.] To propose and evaluate a zero-shot multimodal fusion strategy for species classification based on the integration of visual similarity and textual similarity.

\item[4.] To analyze the conditions and limitations that affect the effectiveness of multimodal models and visual-textual fusion in camera-trap workflows.

\item[5.] To synthesize the empirical findings and proposed methods into a final contribution that clarifies when and how multimodal models can reduce dependence on task-specific training and annotated camera-trap data.

\end{description}

\subsection{Intended Contributions}

The intended contributions of this thesis are as follows:

\begin{itemize}

\item An empirical evaluation of multimodal foundation models for multiple stages of the camera-trap workflow, covering empty-image filtering, species classification, and behavior identification.

\item An analysis of multimodal model behavior under low-supervision scenarios, including zero-shot and few-shot settings and, for behavior identification, the use of image sequences.

\item A zero-shot multimodal fusion strategy for species classification that combines CLIP-based visual similarity with language-model-assisted textual similarity.

\item Evidence on the conditions under which multimodal models and visual-textual fusion can improve camera-trap analysis, as well as the cases in which these approaches remain limited.

\item A methodological basis for reducing dependence on large annotated datasets in camera-trap workflows while preserving practical relevance for biodiversity monitoring.

\end{itemize}

\section{Organization of the Document}

The remainder of this thesis is organized as follows. Chapter 2 presents the fundamental concepts required to understand the work, including zero-shot and few-shot learning, multimodal large language models, image preprocessing, training techniques, and evaluation metrics. Chapter 3 reviews related work on automated camera-trap analysis, focusing on empty-image filtering, species recognition, behavior recognition, multimodal learning, and fusion-based approaches. Chapter 4 presents the first empirical study, which evaluates multimodal models in three camera-trap tasks: empty-image filtering, species classification, and behavior identification. Chapter 5 presents the proposed multimodal fusion strategy for zero-shot species classification and evaluates it across different datasets and image settings. Finally, Chapter 6 summarizes the thesis contributions and outlines the next steps and schedule for completing the research.